% =====================================================================
% Appendix: FEMM current-flow model of the MHD pump (incomplete work)
% Source files: Part 4 Project/FEMM (MHD_pump.lua and report/)
% =====================================================================
\chapter{FEMM Model of the MHD Pump}
\label{app:femm}

A two-dimensional FEMM~4.2 current-flow model of the DC conduction pump was
started in this project to find how the drive current spreads through the
Galinstan channel, and from that the pumping force, the back-EMF and the pump
efficiency. The model was not completed or validated against measurements of
the pump, so it is not presented in the body of the report. This appendix
summarises the work to date as a starting point for the pump redesign
recommended in \cref{sec:disc-further-work}. All cases are generated by one
parametric Lua script, \texttt{MHD\_pump.lua}, held in the \texttt{FEMM} folder
of the project files.
\issue{Confirm that the modelled channel ($30 \times \SI{5}{\milli\metre}$),
\SI{20}{\milli\metre} electrodes and \SI{0.45}{\tesla} field match the built pump;
if not, describe the model as a representative geometry.}

\section{Model}

The model is drawn in the plane of the flow and the current, with the channel
depth as FEMM's out-of-plane depth (\cref{fig:femm-geometry}). The channel is
\SI{30}{\milli\metre} wide and \SI{5}{\milli\metre} deep, with copper electrodes
\SI{20}{\milli\metre} long, and extends \SI{120}{\milli\metre} beyond each
electrode so that the current fringing past the electrode ends is fully
captured. A current of \SI{18}{\ampere} is driven between the electrodes. FEMM's
current-flow solver has no magnetic field, so the force is found by integrating
$\mathbf{J} \times \mathbf{B}$ over a prescribed uniform field of
\SI{0.45}{\tesla}, and the flow is included by adding the motional EMF
$\mathbf{u} \times \mathbf{B}$ for plug flow of up to
\SI{5}{\milli\litre\per\second}. The resistance changed by less than
\SI{0.1}{\percent} when the mesh was refined twice.

\begin{figure}[htbp]
  \centering
  \includegraphics[width=0.6\linewidth]{femm_geometry}
  \caption{Geometry and boundary conditions of the FEMM pump model.}
  \label{fig:femm-geometry}
\end{figure}

\section{Results to Date}

At \SI{18}{\ampere}, the electrode-to-electrode resistance is
\SI{69.6}{\micro\ohm}, and about \SI{40}{\percent} of the current fringes past
the electrode ends, outside the field, where it produces no force
(\cref{fig:femm-jmap}). The model was reduced to a lumped circuit, in which the
Galinstan forms a parasitic path between the electrodes and a fringe path beyond
them, and this circuit reproduced the FEMM force and back-EMF to within
\SI{0.4}{\percent}. From it, the efficiency of the pump has a single peak
(\cref{fig:femm-efficiency}), whose height depends only on the electrode aspect
ratio and the series resistance. \Cref{tab:femm-results} summarises the main
design results.

\begin{figure}[htbp]
  \centering
  \includegraphics[width=0.6\linewidth]{femm_J_map}
  \caption{Current density in the pump channel at \SI{18}{\ampere}, no flow.}
  \label{fig:femm-jmap}
  \floatnote{Streamlines \SI{1}{\ampere} apart.}
\end{figure}

\begin{figure}[htbp]
  \centering
  \includegraphics[width=\linewidth]{femm_efficiency}
  \caption{Modelled pump efficiency against flow and current.}
  \label{fig:femm-efficiency}
  \floatnote{\SI{20}{\milli\metre} electrodes, \SI{0.45}{\tesla}; terminal
    efficiency, excluding leads and contacts.}
\end{figure}

\begin{table}[htbp]
  \centering
  \caption{Main results of the FEMM pump model.}
  \label{tab:femm-results}
  \small
  \begin{tabular}{@{}lc@{}}
    \toprule
    Quantity & Value \\
    \midrule
    Electrode-to-electrode resistance & \SI{69.6}{\micro\ohm} \\
    Current inside the field & \SI{60}{\percent} \\
    Zero-flow pressure rise at \SI{18}{\ampere} & \SI{0.98}{\kilo\pascal} \\
    Efficiency at \SI{5}{\milli\litre\per\second}, \SI{18}{\ampere} & \SI{14.5}{\percent} \\
    Peak efficiency, \SI{20}{\milli\metre} electrodes & \SI{18.8}{\percent} \\
    Peak efficiency, \SI{60}{\milli\metre} electrodes & \SI{36.9}{\percent} \\
    Peak efficiency, \SI{20}{\milli\metre} electrodes,
      \SI{50}{\micro\ohm} series resistance & \SI{12.7}{\percent} \\
    \bottomrule
  \end{tabular}
  \floatnote{Terminal efficiency $\Delta p\,Q/(VI)$; the field covers the
    electrodes.}
\end{table}

The design studies suggest three ways to improve the pump. Lengthening the
electrodes, with the field lengthened to match, brings more of the current into
the field and roughly doubles the peak efficiency at \SI{60}{\milli\metre}.
Widening the channel adds flow area but not pressure unless the field is
lengthened with it. A shallower channel gives more pressure per ampere, but at
the cost of higher viscous loss.

The model also points to a likely cause of the low measured efficiency. The
measured pump power of \SI{0.85}{\watt} at \SI{18}{\ampere} corresponds to a
terminal resistance of about \SI{2.6}{\milli\ohm}, nearly 40 times the
\SI{69.6}{\micro\ohm} of the modelled channel and electrodes. Most of the input
power is therefore likely dissipated in the leads, joints and
electrode--Galinstan contacts, which the model does not include. Because the
pump itself is only tens of microohms, even \SI{50}{\micro\ohm} of series
resistance reduces its peak efficiency by about a third
(\cref{tab:femm-results}), so the supply should sit next to the pump on short,
heavy busbars.

\section{Work Remaining}

The model is two-dimensional, so the top and bottom walls, Hartmann layers and
the real velocity profile are not captured, and it assumes a uniform field with
sharp edges, nominal room-temperature conductivities and ideal contacts. To be
used for design, it would need to be validated against measurements of the
pump's pressure--flow curve and terminal voltage, extended to include the
contact and lead resistances and the fringing of the magnet field, and coupled
to the pressure drop of the cooling loop to find the operating point.
